\subsection{\texorpdfstring{第~\ref{sec:orientation-and-volume-forms}~节}{定向与体积形式}}\label{subsec:solutions-orientation-and-volume-forms}

\begin{description}[font=\normalfont,leftmargin=2em,labelwidth=1.5em]
\RGsolution{ex:forms-and-orientation-on-projective-space}{forms-and-orientation-on-projective-space}
\begin{enumerate}[label=(\alph*),ref=\alph*,leftmargin=*]
\item 若 $\omega\in\Omega^k(S^n)$ 满足 $\omega=\pi^*\theta$, 其中 $\theta\in\Omega^k(\mathbb{R}P^n)$, 则
\[
  f^*\omega=f^*(\pi^*\theta)=(\pi\circ f)^*\theta=\pi^*\theta=\omega.
\]
反之, 设 $\omega\in\Omega^k(S^n)$ 满足 $f^*\omega=\omega$. 对每个使 $\pi|_U$ 为微分同胚的开集 $U\subset S^n$, 定义 $\theta_U\in\Omega^k(\pi(U))$ 为 $\theta_U:=(\pi|_U^{-1})^*\omega$. 现在证明, 若 $p\in\pi(U)\cap\pi(V)$, 其中 $V\subset S^n$ 是另一个使 $\pi|_V$ 为微分同胚的开集, 则 $(\theta_U)_p=(\theta_V)_p$. 实际上, 有 $\pi^{-1}(p)=\{p_1,p_2\}$, 其中 $p_1,p_2\in S^2$ 且 $p_2=f(p_1)$. 不妨设 $p_1\in U$. 则或者 $p_1\in V$, 或者 $p_2\in V$. 若 $p_1\in V$, 则 $\pi|_U$ 和 $\pi|_V$ 在 $p_1$ 的某个邻域内相同, 因此 $\theta_U$ 和 $\theta_V$ 在这个邻域的像 (包含 $p$) 上也相同. 若 $p_2\in V$, 则在 $p_1$ 的某个邻域中有 $\pi|_U=\pi|_V\circ f$, 因而在该邻域上有
\[
  \pi|_U^{-1}=f^{-1}\circ\pi|_V^{-1}=f\circ\pi|_V^{-1}.
\]
因此, 在这个邻域的像上 (特别是在 $p$ 处) 有
\[
  \theta_U=(f\circ\pi|_V^{-1})^*\omega=(\pi|_V^{-1})^* f^*\omega
  =(\pi|_V^{-1})^*\omega=\theta_V.
\]
综上, 在 $\pi(U)\cap\pi(V)$ 上有 $\theta_U=\theta_V$, 从而存在 $\theta\in\Omega^k(\mathbb{R}P^n)$, 使得对每个 $U$ 有 $\theta_U=\theta|_{\pi(U)}$. 最后, 由于对每个 $U$ 有 $\omega|_U=\pi|_U^*\theta_U$, 所以 $\omega=\pi^*\theta$.
\item\label{sol:sphere-volume-form-and-projective-orientability} 容易验证
\[
  \omega=\sum_{i=1}^{n+1}(-1)^{i+1}x^i\,\mathrm{d}x^1\wedge\cdots\wedge
  \mathrm{d}x^{i-1}\wedge\mathrm{d}x^{i+1}\wedge\cdots\wedge\mathrm{d}x^{n+1}
\]
是 $S^n$ 的一个体积形式 [参见 Exercise~\ref{ex:stokes-theorem}(\ref{ex:sphere-vector-field-and-antipodal-homotopy})]. 实际上, 若 $v_1,\ldots,v_n\in\mathbb{R}^{n+1}$ 是 $n$ 个在 $x$ 处与 $S^n$ 相切的线性无关向量, 则
\begin{align*}
  \omega(v_1,\ldots,v_n)
  &=\sum_{i=1}^{n+1}(-1)^{i+1}x^i
  \begin{vmatrix}
  v_1^1&\ldots&v_1^{i-1}&v_1^{i+1}&\ldots&v_1^{n+1}\\
  \vdots&&\vdots&\vdots&&\vdots\\
  v_n^1&\ldots&v_n^{i-1}&v_n^{i+1}&\ldots&v_n^{n+1}
  \end{vmatrix}\\
  &=\begin{vmatrix}
  x^1&\ldots&x^{n+1}\\
  v_1^1&\ldots&v_1^{n+1}\\
  \vdots&&\vdots\\
  v_n^1&\ldots&v_n^{n+1}
  \end{vmatrix}\ne0,
\end{align*}
因为 $\{x,v_1,\ldots,v_n\}$ 是 $\mathbb{R}^{n+1}$ 的一组基. 由于
\[
  f^*\omega=(-1)^{n+1}\omega,
\]
当 $n$ 为奇数时, 有 $\omega=\pi^*\theta$, 其中 $\theta\in\Omega^n(\mathbb{R}P^n)$ 必为体积形式 (因为 $\pi$ 是局部微分同胚). 因此, 对奇数 $n$, $\mathbb{R}P^n$ 可定向. 现在设 $n$ 为偶数. 若 $\mathbb{R}P^n$ 可定向, 就存在体积形式 $\theta\in\Omega^n(\mathbb{R}P^n)$. 此时 $\widetilde\omega:=\pi^*\theta$ 是 $S^n$ 上满足 $f^*\widetilde\omega=\widetilde\omega$ 的体积形式. 由于 $\omega$ 也是 $S^n$ 的体积形式, 有 $\widetilde\omega=g\omega$, 其中 $g\in C^\infty(S^n)$ 处处不为零. 但因为 $f^*\omega=-\omega$, 该函数必须满足 $g\circ f=-g$, 即它在对径点处必须取相反的符号, 因而必在某点为零. 因此, 对偶数 $n$, $\mathbb{R}P^n$ 不可定向. 最后, 设 $n$ 为奇数, $U\subset S^n$ 是开半球. 则 $\pi|_U$ 为微分同胚, 且 $\mathbb{R}P^n\setminus\pi(U)$ 的测度为零, 因此
\begin{align*}
  \int_{S^n}\pi^*\theta
  &=\int_U\pi^*\theta+\int_{f(U)}\pi^*\theta
  =\int_U\pi^*\theta+\int_U f^*\pi^*\theta\\
  &=2\int_U\pi^*\theta=2\int_{\mathbb{R}P^n}\theta.
\end{align*}
\item 考虑由 (\ref{sol:sphere-volume-form-and-projective-orientability}) 中的体积形式 $\omega$ 定义的 $S^n$ 上的定向. 若 $\widetilde{\mathbb{R}P^n}$ 是 $\mathbb{R}P^n$ 的可定向二重覆盖 [参见第~\ref{chap:differentiable-manifolds}~章的 Exercise~\ref{ex:orientability}(\ref{ex:orientable-double-covering})], 定义 $g:S^n\to\widetilde{\mathbb{R}P^n}$ 为
\[
  g(p)=(\pi(p),[(\mathrm{d}\pi)_p v_1,\ldots,(\mathrm{d}\pi)_p v_n]),
\]
其中 $\{v_1,\ldots,v_n\}$ 是 $T_pS^n$ 的一组正定向基. 对偶数 $n$, 有 $f^*\omega=-\omega$, 即 $f$ 反转定向. 据此容易证明 $g$ 是满足 $\widetilde\pi\circ g=\pi$ 的微分同胚, 其中 $\widetilde\pi:\widetilde{\mathbb{R}P^n}\to\mathbb{R}P^n$ 是自然投影.
\end{enumerate}

\RGsolution{ex:brouwer-fixed-point-theorem}{brouwer-fixed-point-theorem}
\begin{enumerate}[label=(\alph*),ref=\alph*,leftmargin=*]
\item\label{sol:no-smooth-retraction-onto-boundary} 设 $\omega\in\Omega^{n-1}(\partial M)$ 是 $\partial M$ 上与诱导定向相容的体积形式. 由于 $\partial M$ 是紧的 (因为 $M$ 是紧的), 有
\[
  \int_{\partial M}\omega>0.
\]
若 $f$ 存在, 则它满足 $f\circ i=\operatorname{id}$, 其中 $i:\partial M\to M$ 是包含映射. 此时由 Stokes 定理得
\begin{align*}
  \int_{\partial M}\omega
  &=\int_{\partial M}(f\circ i)^*\omega=\int_{\partial M}i^* f^*\omega
  =\int_M\mathrm{d}f^*\omega\\
  &=\int_M f^*\mathrm{d}\omega=\int_M f^*0=0
\end{align*}
(因为 $\mathrm{d}\omega$ 是 $(n-1)$ 维流形 $\partial M$ 上的 $n$-形式, 故 $\mathrm{d}\omega=0$). 因此 $f$ 不可能存在.
\item 假设存在没有不动点的可微映射 $g:B\to B$. 那么对每个 $x\in B$, 存在唯一一条以 $g(x)$ 为起点并经过 $x$ 的射线 $r_x$, 而 $r_x\setminus\{g(x)\}$ 显然在唯一一点 $f(x)$ 与 $\partial B$ 相交 (参见 Fig.~\ref{fig:brouwer-retraction-from-ball}). 映射 $f:B\to\partial B$ 将是可微的, 因为
\[
  f(x)=g(x)+t(x)\frac{x-g(x)}{\lVert x-g(x)\rVert},
\]
其中 $t(x)$ 是方程
\[
  \lVert f(x)\rVert^2=1
  \iff t^2+2t\frac{\langle x,g(x)\rangle-\lVert g(x)\rVert^2}{\lVert x-g(x)\rVert}
           +\lVert g(x)\rVert^2=1
\]
的唯一正根. 此外, 有 $f|_{\partial B}=\operatorname{id}$. 由 (\ref{sol:no-smooth-retraction-onto-boundary}) 知 $f$ 不可能存在, 故 $g$ 也不可能存在.
\end{enumerate}
\end{description}

\RGfigure[.42\textwidth]{brouwer-retraction-from-ball}{映射 $f:B\to\partial B$ 的构造}
